\begin{table}
\centering
\caption{\label{tab:tab_metric_breakdown}Indicators and metrics by dimension.}
\centering
\fontsize{10}{12}\selectfont
\begin{threeparttable}
\begin{tabular}[t]{lrlrrrr}
\toprule
\multicolumn{1}{c}{ } & \multicolumn{2}{c}{Indicator} & \multicolumn{4}{c}{Metric} \\
\cmidrule(l{3pt}r{3pt}){2-3} \cmidrule(l{3pt}r{3pt}){4-7}
Dimension & N & Represented & N & County & Trend & County Trend\\
\midrule
Economics & 12 & 8 (67\%) & 11 & 9 & 10 & 9\\
Environment & 13 & 10 (77\%) & 16 & 10 & 14 & 9\\
Health & 13 & 11 (85\%) & 14 & 11 & 12 & 10\\
Production & 12 & 10 (83\%) & 12 & 5 & 11 & 4\\
Social & 20 & 8 (40\%) & 8 & 7 & 7 & 6\\
\addlinespace
\textbf{Total} & \textbf{70} & \textbf{47 (67\%)} & \textbf{61} & \textbf{42} & \textbf{54} & \textbf{38}\\
\bottomrule
\end{tabular}
\begin{tablenotes}
\item \textit{Note:} Representation is the percentage of indicators represented by at least one metric. County metrics are available at or below the county level. Trend metrics have at least two time points for at least one state or county; County Trend metrics meet both requirements for trend analyses.
\end{tablenotes}
\end{threeparttable}
\end{table}
